\documentclass[10pt,twocolumn]{article}
\usepackage[margin=0.72in,columnsep=0.24in]{geometry}
\usepackage[T1]{fontenc}
\usepackage[utf8]{inputenc}
\usepackage{lmodern,microtype}
\usepackage{amsmath,amssymb,booktabs,tabularx,array,xcolor}
\usepackage{enumitem}
\usepackage{balance}
\usepackage[hidelinks]{hyperref}
\definecolor{authority}{HTML}{0B6E75}
\setlist{nosep,leftmargin=*}
\setlength{\emergencystretch}{2em}
\newcommand{\system}{\textsc{AtMem}}
\title{\textbf{When the Write Lands but the Acknowledgement Is Lost:\\A Reproducible Memory-Retry Test in OpenClaw}}
\author{Javad Taghia\\AtMem.Ai Lab\\\url{https://atmem.ai}}
\date{September 27, 2026}
\begin{document}
\raggedbottom
\twocolumn[
\begin{@twocolumnfalse}
\maketitle
\begin{abstract}
A memory write can succeed even when its caller receives a timeout. Retrying
then raises a question that recall accuracy alone cannot answer: does the
agent read one complete original record, a duplicate, or conflicting state?
We present a small, reproducible fault-injection study of the AtMem
\texttt{memory\_remember} integration with OpenClaw. A transport proxy observes
a successful backend response, withholds it, and kills the backend process;
the caller receives a real timeout and retries. The final qualification uses
five isolated bridge trials and three real OpenClaw gateway trials, plus one
gateway no-fault control. In every faulted qualification trial, retries return
the original record identifier with status \texttt{already\_stored}; the final
store contains one active record with the complete original fact and a valid
audit chain. Gateway trials use deterministic local model responses, not
paid inference or an autonomous learned retry policy. We retain unsuccessful
CLI setup attempts separately. These results establish recovery of one exact
synthetic fact after a completed write and lost response in the tested
configuration. They do not establish transaction-interruption safety,
paraphrase deduplication, or a general exactly-once guarantee.
\end{abstract}
\vspace{0.6em}
\noindent\textbf{Keywords:} persistent memory, OpenClaw, acknowledgement loss,
tool timeout, retry, fault injection, reproducibility
\vspace{1em}
\end{@twocolumnfalse}
]

\section{The question}
Consider an agent instructed to remember a preference. Its tool sends a write
to a persistent memory service. The service commits the record, but the reply
never reaches the agent. The agent sees a timeout, retries, and reads back the
memory. Returning the right words is necessary but insufficient: the retry
could have created two records, changed the source linkage, or overwritten
metadata. A useful test therefore inspects both the agent-visible tool result
and the persisted state.

This note follows the broader integrity and continuity evaluation in
\emph{Beyond Recall Accuracy}~\cite{taghia2026beyond}. Its question is narrower:
\emph{what state does an OpenClaw agent read after a successful memory write
whose acknowledgement was lost?} The experiment tests an actual memory write,
rather than an external document publication or retail exchange.

\paragraph{Two different failure windows.}
A partially completed transaction and a completed write with a lost response
are distinct. Here the fault is injected only after the backend has produced a
successful response containing a new record. This is a post-write failure
window. Killing the backend before the response is delivered does not test
every point inside its storage transaction or derived-index updates.

\section{System and test boundary}
The tested path is OpenClaw's registered \texttt{memory\_remember} tool,
the production TypeScript AtMem bridge, and the real Python MCP backend
with a file-backed SQLite memory store. Only the transport proxy and test
driver are experimental. Product write, duplicate handling, record read, and
audit-verification implementations are unchanged.

Qualification uses a clean exported source snapshot at commit
\nolinkurl{da046dc59fd41fdd77ed407d8f196cfc434b478c}.
The bridge package declares version 2.3.7. The installed OpenClaw runtime
reports 2026.9.1 (\texttt{ad6fe23}). This is a source-snapshot qualification,
not certification of a newly released wheel. Runtime versions and source
hashes are recorded in the accompanying artifact manifest.

Each trial starts with a new directory, memory database, and synthetic
subject. Gateway trials additionally use an isolated OpenClaw state directory,
workspace, plugin installation, loopback model endpoint, and gateway. They do
not modify the user's normal agent memory or send messages to external
channels. No paid model calls are made.

\subsection{Two execution levels}
\textbf{Bridge trials} register the actual bridge plugin against a minimal
OpenClaw API harness. The harness supplies a synthetic authenticated-message
binding and invokes the registered tool. It retries first with the original
tool-call identifier and then with a new identifier. These trials directly
inspect the complete record before and after retry.

\textbf{Gateway trials} run the installed OpenClaw gateway and ordinary agent
command. A local OpenAI-compatible endpoint supplies a fixed sequence:
write, retry with a new tool-call identifier, read the returned record, finish.
OpenClaw executes the real tools and passes their actual results back to the
endpoint. Its normal lifecycle hooks stage the source message. The endpoint
chooses the read identifier from the observed tool result, not from the
evaluator's private copy of the withheld response.

This second level establishes host integration, not language-model reasoning
quality. The retry policy is scripted. A learned model might retry differently,
change the wording, or decline to retry; those behaviors are outside this study.

\section{Protocol and assertions}
The synthetic user request is to remember a preference for cobalt blue
notebooks. The proposed memory is exactly:
\begin{quote}\small
User prefers cobalt blue notebooks.
\end{quote}
The replaceable slot is \texttt{notebook\_preference}. Source type is
\texttt{user\_message}. Each trial follows this sequence:
\begin{enumerate}
\item Stage the source message through the normal host hook, or explicitly
through the real staging tool in the bridge harness.
\item Invoke \texttt{memory\_remember} with the synthetic fact and slot.
\item The proxy waits for a successful MCP reply containing a newly created
record. It saves the request and response as evaluator evidence, never forwards
that reply, and sends \texttt{SIGKILL} to the backend process.
\item Keep the transport open so the real bridge request deadline expires
after 2,000\,ms. The caller receives a timeout rather than a success response.
\item Reconnect to a fresh backend process using the same database and retry
the identical fact. Bridge trials exercise both unchanged and new call IDs;
gateway trials exercise a new call ID.
\item Read the record through \texttt{memory\_get}, enumerate all records
including inactive ones, and verify retained audit evidence.
\end{enumerate}

The gateway endpoint waits 600\,ms before returning each scripted response,
allowing the bridge's 250\,ms idle-close interval to dispose of the faulted
transport. Bridge trials explicitly stop their registered service before
reconnecting. Immediate retry on a still-open faulted connection is not
qualified by this schedule. The operating system and gateway remain alive;
the process killed is the memory backend.

\paragraph{Success criteria.}
A faulted trial must show a committed record in the withheld reply, an actual
caller-visible timeout, the original identifier in the successful retry,
\texttt{stored=true}, \texttt{status=already\_stored}, one final active record,
the exact original fact on read-back, one \texttt{memory.record\_created}
audit event, and a valid audit chain. Bridge qualification additionally
requires equality of the entire record list before and after both retries.
Gateway qualification additionally checks SQLite integrity and foreign-key
consistency. Normal host termination alone is not a pass.

\section{Results}
Table~\ref{tab:results} separates execution levels and the control. Counts
describe repeated deterministic engineering trials, not sampled production
failure probabilities. All final qualification trials use the same fact and
fault boundary, with independent temporary stores.

\begin{table}[ht]
\centering\small
\caption{Final qualification. Each successful run ends with one record.
The two retries in a bridge trial are sequential, not independent trials.}
\label{tab:results}
\begin{tabularx}{\columnwidth}{@{}Xrr@{}}
\toprule
Condition & Runs & Passes\\
\midrule
Bridge: lost response, two retry IDs & 5 & 5\\
OpenClaw gateway: lost response & 3 & 3\\
OpenClaw gateway: no-fault control & 1 & 1\\
\bottomrule
\end{tabularx}
\end{table}

\subsection{What the agent reads back}
The lost-response gateway sequence exposes a timeout on the first call.
The retry returns the original record ID, \texttt{stored=true}, and
\texttt{already\_stored}. The subsequent \texttt{memory\_get} returns
``User prefers cobalt blue notebooks.'' Its provenance identifies
\texttt{typed\_session\_handoff} and a host-asserted semantic interpretation.

Inspection of the final stored record shows the original source turn, original
episode linkage, \texttt{trusted\_user} trust tier, active status, the same
fact key, generation zero, and no superseding record. Including inactive
records does not reveal a hidden duplicate. In the bridge trials, the complete
record dictionaries before and after both retries are identical, including
creation timestamp and metadata.

The no-fault gateway control receives success on the first write and
\texttt{already\_stored} on its deliberate repeat. This confirms that the
scripted retry and read-back also work without the injected failure.

\subsection{Record identity and retry events}
The audit trail retains the original creation and duplicate-recognition
events. There is one record-creation event per trial, but retries are still
invocations and are still observable. Bridge trials record two semantic
duplicate events; gateway trials record one. We therefore describe the result
as \emph{one preserved canonical record under exact-content retry}, rather
than claiming that every side effect occurred exactly once.

\subsection{Development and setup outcomes}
Before final qualification, six bridge development trials passed. Two embedded
OpenClaw CLI runs failed before reaching the fault boundary: no current typed
user-message binding was available, so the write was rejected and the store
remained empty. These are failed setup attempts, not successful crash tests
or memory corruption. Changing to an isolated gateway startup allowed source
staging and the subsequent development gateway trial passed.

The final qualification was then run from the clean source snapshot. The
artifact index retains the development outcomes and marks their different
provenance. The two failed embedded runs remain an integration limitation;
this experiment does not establish their root cause or claim a product fix.

\balance
\section{Interpretation and limitations}
The direct answer to the motivating question is: \emph{in these tests, the
agent reads the same complete original memory, and its retry is reported as
already stored}. The result is stronger than checking that recall returns the
expected phrase because it also checks identity, record multiplicity,
metadata preservation, and audit evidence.

It is nevertheless narrow. A new call identifier still used identical fact
bytes and the same slot and source binding. This demonstrates the observed
duplicate handling; it does not establish an API-wide idempotency-key contract.
Changing the fact, paraphrasing it, changing the source, or racing concurrent
writers could produce different behavior. No baseline or competing memory
system is evaluated, and no superiority claim follows.

The backend was killed after emitting a successful response, not during
individual SQL statements. No whole-machine crash, sudden power loss, torn
storage write, filesystem failure, or distributed replica failure was injected.
SQLite integrity checks and valid audit hashes do not prove those untested
properties. Nor do they independently establish the truth of a user's fact.

The host run used an explicitly enabled legacy direct-memory tool profile,
not every managed control-plane configuration or arbitrary OpenClaw tool.
The deterministic endpoint tests actual tool dispatch and result handling,
but not autonomous recovery planning. The deliberately paced retry follows
transport reconnection; alternate retry timing remains untested.

\section{Reproduction and artifacts}
The public supplement contains the LaTeX source, PDF, probe programs,
clean-source snapshot, per-run JSON evidence, a machine-readable summary,
runtime and source hashes, and reproduction instructions. It is published
alongside the earlier paper in the
\href{https://huggingface.co/datasets/atmem/memory-integrity-continuity/tree/main/openclaw-lost-ack-20260927}{Hugging Face evidence dataset}.
The source snapshot and original backend response make the fault boundary
inspectable. The saved response is evaluator-only evidence; the agent receives
the timeout and must retry to obtain its record identifier.

No existing results or earlier paper are replaced. The supplement is a
functional qualification note and should not be interpreted as a new
leaderboard score or a general reliability rate.

\begin{thebibliography}{1}
\bibitem{taghia2026beyond}
Javad Taghia.
\newblock Beyond Recall Accuracy: Evaluating Integrity and Crash Continuity
in Persistent Memory for Tool-Using Language Agents.
\newblock AtMem.Ai Lab technical manuscript, September 26, 2026.
\newblock \href{https://huggingface.co/datasets/atmem/memory-integrity-continuity/blob/main/paper/atmem-memory-integrity-and-continuity.pdf}{Manuscript and accompanying evidence on Hugging Face}.
\end{thebibliography}
\end{document}
